% Source-aligned English for printed p.208.
\clearpage
\SourcePage{208}

\LayerTarget{tohoku:proof:lemma-5.5.2:p1}
\noindent of injective objects in $\mathbf K({\mathbf C}_1)$ given in 2.4, $\mathrm C$ is
acyclic in dimensions $>0$ and ``decomposes''; moreover, the $\mathrm C^i$ and
$Z^i(\mathrm C)$ are injective. It follows first that $\overline F(\mathrm C)$ is a resolution
of $F(Z^0(\mathrm C))$, whence
${\mathbf R}^nG(\overline F(\mathrm C))=R^nG(F(Z^0(\mathrm C)))$; and since $Z^0(\mathrm C)$ is injective, the
second member vanishes by the hypothesis on $F$. We may therefore apply
Theorem 2.4.1 to the composite functor $(H^0\overline G)\overline F$. This gives
the spectral functor in the statement, taking account of
$R^q\overline F=\overline{R^qF}$ (a relation following immediately from the
definitions).

\LayerTarget{tohoku:corollary:lemma-5.5.2}
\textsc{Corollary.} \textit{Under the preceding conditions, let $\mathrm C$
be a resolution of an object $A\in {\mathbf C}_1$. Then there exists a cohomological
spectral sequence abutting to the graded object $(R^n(GF)(A))$, whose initial
term is given in \textup{(5.5.2)}.}

\LayerTarget{tohoku:proof:corollary-lemma-5.5.2}
Indeed, here ${\mathbf R}^n(GF)(\mathrm C)=R^n(GF)(A)$.

\LayerTarget{tohoku:definition:g-cover}
Return to the space $X$ with its operator group $G$. Let
${\mathbf U}=(U_i)_{i\in I}$ be a covering of $X$, with $G$ acting on the index
set $I$ so that $g. U_i=U_{g. i}$ for every $i\in I$. We shall call
${\mathbf U}$ a $G$-\textit{covering}. If $P$ is an abelian presheaf on $X$ on which $G$
acts, then the complex $C({\mathbf U},P)$ of cochains of ${\mathbf U}$ with
coefficients in $P$ may obviously be regarded as a complex of $G$-groups. We
put

\LayerTarget{tohoku:equation:5.5.3}
\[
\tag{5.5.3}
H^n({\mathbf U};G,P)={\mathbf R}^n\Gamma^G(C({\mathbf U},P)),
\]
\LayerTarget{tohoku:construction:equivariant-cech-bicomplex}
where the second member may therefore also be written as the group
$H^n(C(G,C({\mathbf U},P)))$. Here $C(G,C({\mathbf U},P))$ is the bicomplex
formed by the cochains of $G$ with values in the complex $C({\mathbf U},P)$, a
bicomplex which may be denoted $C({\mathbf U};G,P)$:

\LayerTarget{tohoku:equation:5.5.3-bis}
\[
\tag{5.5.3 bis}
H^n({\mathbf U};G,P)=H^n(C({\mathbf U};G,P)).
\]

\LayerTarget{tohoku:construction:open-g-cover-spectral-functor}
Suppose first that the $G$-covering ${\mathbf U}$ is \textit{open}. We apply Lemma 5.5.1
with the categories ${\mathbf C}^{X(G)}$, ${\mathbf C}^G$ (the category of $G$-modules), and ${\mathbf C}$
(the category of abelian groups), the functors $\Gamma_X$ and $\Gamma^G$, and
finally $\mathbf K(A)=C({\mathbf U},A)$. The functor $\Gamma_X$ sends injective
objects to $\Gamma^G$-acyclic objects (the corollary to Proposition 5.1.3), and
conditions (i) and (ii) of the lemma are satisfied: (i) because ${\mathbf U}$ is
open, and (ii) by what was said in 3.8 (second paragraph). It remains to
identify $R^q\mathbf K(A)$. The calculation is immediate (and was already made
in 3.8): one finds $R^q\mathbf K(A)=C({\mathbf U},H^q(A))$, where $H^q(A)$
denotes the presheaf $H^q(A)(V)=H^q(V,A)$. We have therefore proved:

\LayerTarget{tohoku:theorem:5.5.3}
\textsc{Theorem 5.5.3.} \textit{Let $X$ be a space equipped with an
operator group $G$, and let ${\mathbf U}=(U_i)$ be an open $G$-covering of $X$.
Then on the category ${\mathbf C}^{X(G)}$ of abelian $G$-sheaves on $X$ there exists a
cohomological spectral functor abutting to the graded functor
$(H^n(X;G,A))$, whose initial term is}

\LayerTarget{tohoku:equation:5.5.4}
\[
\tag{5.5.4}
E_2^{p,q}(A)=H^p({\mathbf U};G,H^q(A))
\]
\textit{(where $H^q(A)$ is the presheaf $H^q(A)(V)=H^q(V,A)$ on $X$).}

\LayerTarget{tohoku:map:5.5.5}
As usual, one deduces the edge homomorphisms
\[
\tag{5.5.5}
H^n(U;G,A)\to H^n(X;G,A)
\]
\LayerTarget{tohoku:conclusion:theorem-5.5.3:p1}
and a five-term exact sequence showing, among other things, that \textit{the preceding homomorphism}

\TohokuContinuation{209}
